\begin{proof}[Proof of \cref{obj:lem:ordered-mass}]
\begin{enumerate}
\item \emph{Uniform coefficient.}
Set
\[
m:=\lceil\beta\rceil-1,
\qquad
a_{\mathrm{tpl}}:=c_f\eta^d,
\]
using the template of \cref{obj:synth_1,obj:synth_2,obj:def:template}. Its polynomial coordinates and node set are
\[
\mathcal I_m=\{\alpha\in\mathbb N_0^d:|\alpha|\le m\},
\qquad
U(z)=(z^\alpha)_{\alpha\in\mathcal I_m},
\qquad
J=(m+1)^d,
\]
\[
\{v_\ell:1\le\ell\le J\}
=
\left\{
\left(\frac{\iota_i+1}{m+2}\right)_{i=1}^d:
\iota\in\{0,\ldots,m\}^d
\right\}.
\]
We first establish positivity of the template width. Define
\[
\mathcal B_m:=
\left\{
b\in\mathbb R^{\mathcal I_m}:
\sum_{\alpha\in\mathcal I_m}b_\alpha^2=1
\right\},
\qquad
\pi_b(z):=\sum_{\alpha\in\mathcal I_m}b_\alpha z^\alpha
\quad(z\in\mathbb R^d).
\]
For \(b\in\mathcal B_m\), at least one coefficient is nonzero, and
\[
\pi_b\not\equiv0,
\qquad
\deg_{z_i}\pi_b\le m
\quad(1\le i\le d).
\]
The \(m+1\) node values in each coordinate are distinct. If \(\pi_b\) vanished at every tensor node, fixing the other coordinates at node values and applying the univariate root bound would make it vanish identically in the first coordinate. Repeating this argument in each coordinate would give \(\pi_b\equiv0\). Thus
\[
\exists\,\ell\in\{1,\ldots,J\}:\ \pi_b(v_\ell)\ne0,
\qquad
b^\top G_0b
=
\frac1J\sum_{\ell=1}^J\pi_b(v_\ell)^2>0.
\]
% lean: CausalSmith.Stat.WeakOverlap.templateGram_unit_pos

The set \(\mathcal B_m\) is nonempty, since it contains the coefficient vector of the constant polynomial \(1\). It is closed and bounded: for every \(b\in\mathcal B_m\),
\[
b_\alpha^2\le
\sum_{\alpha'\in\mathcal I_m}b_{\alpha'}^2=1
\quad(\alpha\in\mathcal I_m).
\]
Consequently it is compact. The continuous quadratic form attains its minimum there, and the Rayleigh characterization gives
\[
\lambda_0
=
\min_{b\in\mathcal B_m}b^\top G_0b>0.
\]
% lean: CausalSmith.Stat.WeakOverlap.templateLambda_pos

Every quotient defining \(L_U\) is nonnegative; the indexing set contains distinct cube points because \(d\ge1\). Hence the width rule supplies
\[
L_U\ge0,
\qquad
0<\eta
=
\min\left\{\frac1{2(m+2)},\frac{\lambda_0}{1+L_U}\right\}
\le\frac1{2(m+2)},
\qquad
a_{\mathrm{tpl}}>0.
\]
% lean: CausalSmith.Stat.WeakOverlap.templateLip_nonneg

For \(\gamma'\in\Gamma\), define
\[
\Phi_\Gamma(\gamma')
:=
\frac{\gamma'-1}{\gamma'}\,
C^{-1/(\gamma'-1)}
a_{\mathrm{tpl}}^{\gamma'/(\gamma'-1)}.
\]
Since \(C\ge1\), \(a_{\mathrm{tpl}}>0\), and
\(\gamma'\ge\gamma_{\min}>1\), this function is continuous and strictly positive on the nonempty compact interval \(\Gamma\). Therefore
\[
\kappa_\Gamma:=\min_{\gamma'\in\Gamma}\Phi_\Gamma(\gamma')>0,
\qquad
\kappa_\Gamma\le
\frac{\gamma-1}{\gamma}\,
C^{-1/(\gamma-1)}
a_{\mathrm{tpl}}^{\gamma/(\gamma-1)}
\quad(\gamma\in\Gamma).
\]
This constant depends only on the fixed parameters and interval endpoints.
% lean: CausalSmith.Stat.WeakOverlap.orderedMass_uniform_coefficient

\item \emph{Setwise treated mass.}
Fix \(\gamma\in\Gamma\) and \(P\in\mathcal P_\gamma\), and write
\[
\nu:=P_X.
\]
Then \(\gamma\ge\gamma_{\min}>1\). Let \(E\subseteq\mathcal X\) be measurable and set
\[
u:=\nu(E),
\qquad 0\le u\le1.
\]
If \(u=0\), the claimed lower bound is zero, while the treated probability is nonnegative. Suppose henceforth that \(u>0\), and define
\[
s_E:=\left(\frac uC\right)^{1/(\gamma-1)}.
\]
The inequalities \(0<u/C\le1\) and \(1/(\gamma-1)>0\) give
\[
0<s_E\le1.
\]
% lean: CausalSmith.Stat.WeakOverlap.orderedMass_threshold_mem_unit

Define the weak level-set mass function on all real thresholds by
\[
F_E(t):=\nu\bigl(E\cap\{x:t\le e(x)\}\bigr),
\qquad t\in\mathbb R.
\]
It is nonnegative and nonincreasing:
\[
0\le F_E(t')\le F_E(t)\le u
\qquad(t\le t').
\]
As a version of the conditional treatment probability, the propensity satisfies
\[
0\le e(x)\le1
\quad\text{for }\nu\text{-almost every }x.
\]
Disintegration of the treatment probability therefore gives
\[
P(A=1,X\in E)=\int_E e(x)\,d\nu(x).
\]
Moreover,
\[
e(x)=\int_0^1\mathbf 1\{t\le e(x)\}\,dt
\quad\text{for }\nu\text{-almost every }x.
\]
Tonelli's theorem yields the weak-level-set representation
\[
P(A=1,X\in E)=\int_0^1F_E(t)\,dt.
\]
The propensity is integrable on \(E\), and the bounded monotone level-set functions are integrable on the threshold intervals used below.
% lean: CausalSmith.Stat.WeakOverlap.orderedMass_treated_real_integral

For comparison, define
\[
F_E^{>}(t):=\nu\bigl(E\cap\{x:t<e(x)\}\bigr),
\qquad t\in\mathbb R.
\]
For \(t\in[0,1]\), \cref{obj:ass:global-tail} and monotonicity of measure give
\[
\nu\bigl(E\cap\{e\le t\}\bigr)
\le \nu(\{e\le t\})
\le Ct^{\gamma-1}.
\]
The two threshold sets cover \(E\), so subadditivity implies
\[
u\le F_E^{>}(t)+\nu\bigl(E\cap\{e\le t\}\bigr).
\]
Since \(\{e>t\}\subseteq\{e\ge t\}\), it follows that
\[
u-Ct^{\gamma-1}\le F_E^{>}(t)\le F_E(t)
\qquad(0\le t\le1).
\]
Integrating over the cutoff interval gives
\[
\int_0^{s_E}(u-Ct^{\gamma-1})\,dt
\le
\int_0^{s_E}F_E^{>}(t)\,dt
\le
\int_0^{s_E}F_E(t)\,dt.
\]
% lean: CausalSmith.Stat.WeakOverlap.orderedMass_layercake_piece_lower

The cutoff balances the tail envelope and set mass:
\[
Cs_E^{\gamma-1}
=
C\left(\frac uC\right)^{(1/(\gamma-1))(\gamma-1)}
=u.
\]
Also \(s_E^\gamma=s_E^{\gamma-1}s_E\). Thus direct integration gives
\[
\begin{aligned}
\int_0^{s_E}(u-Ct^{\gamma-1})\,dt
&=us_E-\frac C\gamma s_E^\gamma\\
&=\frac{\gamma-1}{\gamma}\,us_E\\
&=\frac{\gamma-1}{\gamma}\,
C^{-1/(\gamma-1)}u^{\gamma/(\gamma-1)}.
\end{aligned}
\]
% lean: CausalSmith.Stat.WeakOverlap.orderedMass_integral_at_threshold

Combining these bounds, using \(F_E\ge0\) and \(s_E\le1\), proves
\[
\begin{aligned}
\frac{\gamma-1}{\gamma}\,
C^{-1/(\gamma-1)}u^{\gamma/(\gamma-1)}
&\le \int_0^{s_E}F_E(t)\,dt\\
&\le \int_0^1F_E(t)\,dt\\
&=P(A=1,X\in E).
\end{aligned}
\]
% lean: CausalSmith.Stat.WeakOverlap.orderedMass_setwise_lower

\item \emph{A ranked union and its upper bound.}
Fix \(j\in\mathbb N\) and an integer \(k\) in the stated range, and set
\[
h:=2^{-j}>0,
\qquad
D_\gamma:=\frac{d\gamma}{\gamma-1}.
\]
Use the dyadic partition \(\mathcal Q_h\), whose coordinate intervals are half-open with the global right boundary included. For each cube, write its lower corner and scaled microcells as
\[
a_Q:=\left(\inf_{x\in Q}x_i\right)_{i=1}^d,
\qquad
E_{Q\ell}:=a_Q+hV_\ell
\quad(1\le\ell\le J).
\]
Because the finite collection of microcell masses is nonempty, choose
\[
\ell_{\min}:\mathcal Q_h\longrightarrow\{1,\ldots,J\},
\qquad
q_{Q,\ell_{\min}(Q)}=q_Q
\quad(Q\in\mathcal Q_h).
\]
Define the threshold collection
\[
\mathcal Q_{h,\le k}
:=
\{Q\in\mathcal Q_h:q_Q\le q_{(k)}\}.
\]
The first \(k\) entries of the sorted mass list are all at most \(q_{(k)}\). Sorting preserves the multiplicity of each mass, so
\[
|\mathcal Q_{h,\le k}|\ge k.
\]
Choose a collection and define its selected union by
\[
\mathcal Q_{h,k}\subseteq\mathcal Q_{h,\le k},
\qquad
|\mathcal Q_{h,k}|=k,
\qquad
E_k:=\bigcup_{Q\in\mathcal Q_{h,k}}E_{Q,\ell_{\min}(Q)}.
\]
Subadditivity of the treated probability then gives
\[
\begin{aligned}
P(A=1,X\in E_k)
&\le
\sum_{Q\in\mathcal Q_{h,k}}
q_{Q,\ell_{\min}(Q)}\\
&=
\sum_{Q\in\mathcal Q_{h,k}}q_Q\\
&\le kq_{(k)}.
\end{aligned}
\]
% lean: CausalSmith.Stat.WeakOverlap.orderedMass_ranked_union_upper

\item \emph{The selected union's lower bound.}
The node coordinates and width bound satisfy
\[
\frac1{m+2}\le(v_\ell)_i\le\frac{m+1}{m+2},
\qquad
\frac\eta2\le\frac1{4(m+2)}.
\]
Hence, for every \(z\in V_\ell\) and every coordinate \(i\),
\[
0<(v_\ell)_i-\frac\eta2
\le z_i
\le(v_\ell)_i+\frac\eta2<1.
\]
After scaling, every microcell lies strictly between its cube's coordinate faces:
\[
(a_Q)_i<x_i<(a_Q)_i+h
\quad(x\in E_{Q\ell},\ 1\le i\le d).
\]
In particular, \(E_{Q\ell}\subseteq Q\subseteq\mathcal X\). Distinct dyadic cubes have lower corners differing by at least \(h\) in some coordinate. The strict bounds therefore imply
\[
E_{Q\ell}\cap E_{Q'\ell'}=\varnothing
\quad(Q\ne Q').
\]
% lean: CausalSmith.Stat.WeakOverlap.orderedMass_scaledMicroCell_disjoint_macro

More explicitly, each microcell is the rectangle
\[
E_{Q\ell}
=
\prod_{i=1}^d
\left[
(a_Q)_i+h\left((v_\ell)_i-\frac\eta2\right),
(a_Q)_i+h\left((v_\ell)_i+\frac\eta2\right)
\right].
\]
Each coordinate interval has length \(\eta h\), so
\[
\operatorname{Leb}(E_{Q\ell})=(\eta h)^d.
\]
% lean: CausalSmith.Stat.WeakOverlap.orderedMass_scaledMicroCell_volume

By \cref{obj:ass:density}, integration of the density over each microcell yields
\[
\nu(E_{Q\ell})
=
\int_{E_{Q\ell}}f(x)\,dx
\ge c_f\operatorname{Leb}(E_{Q\ell})
=c_f(\eta h)^d.
\]
The selected microcells are measurable and pairwise disjoint. Thus
\[
\begin{aligned}
\nu(E_k)
&=
\sum_{Q\in\mathcal Q_{h,k}}
\nu(E_{Q,\ell_{\min}(Q)})\\
&\ge k\,c_f(\eta h)^d.
\end{aligned}
\]
% lean: CausalSmith.Stat.WeakOverlap.orderedMass_selectedMicrocells_covariate_lower

The exponent \(\gamma/(\gamma-1)\) and the setwise coefficient are positive. Applying monotonicity of this power and the setwise estimate already established gives
\[
\begin{aligned}
&\frac{\gamma-1}{\gamma}\,
C^{-1/(\gamma-1)}
\bigl[k\,c_f(\eta h)^d\bigr]^{\gamma/(\gamma-1)}\\
&\qquad\le
\frac{\gamma-1}{\gamma}\,
C^{-1/(\gamma-1)}
\nu(E_k)^{\gamma/(\gamma-1)}
\le P(A=1,X\in E_k).
\end{aligned}
\]
% lean: CausalSmith.Stat.WeakOverlap.orderedMass_selectedMicrocells_treated_lower

\item \emph{Power identity and conclusion.}
The factorization and exponent identity are
\[
c_f(\eta h)^d=a_{\mathrm{tpl}}h^d,
\qquad
\frac{\gamma}{\gamma-1}
=
1+\frac1{\gamma-1}.
\]
Using multiplicativity of real powers on positive bases and \(k>0\), we obtain
\[
\begin{aligned}
\frac{\bigl[k\,c_f(\eta h)^d\bigr]^{\gamma/(\gamma-1)}}{k}
&=
a_{\mathrm{tpl}}^{\gamma/(\gamma-1)}
h^{d\gamma/(\gamma-1)}
k^{\gamma/(\gamma-1)-1}\\
&=
a_{\mathrm{tpl}}^{\gamma/(\gamma-1)}
h^{D_\gamma}k^{1/(\gamma-1)}.
\end{aligned}
\]
% lean: CausalSmith.Stat.WeakOverlap.orderedMass_rank_power_identity

Multiplying the uniform coefficient bound by the positive mesh and rank factors therefore gives
\[
\kappa_\Gamma h^{D_\gamma}k^{1/(\gamma-1)}
\le
\frac{
\frac{\gamma-1}{\gamma}\,
C^{-1/(\gamma-1)}
\bigl[k\,c_f(\eta h)^d\bigr]^{\gamma/(\gamma-1)}
}{k}.
\]
The lower and upper bounds on the same selected union give
\[
\frac{\gamma-1}{\gamma}\,
C^{-1/(\gamma-1)}
\bigl[k\,c_f(\eta h)^d\bigr]^{\gamma/(\gamma-1)}
\le P(A=1,X\in E_k)
\le kq_{(k)}.
\]
Dividing by \(k>0\) and combining the two inequalities proves
\[
q_{(k)}
\ge
\kappa_\Gamma h^{D_\gamma}k^{1/(\gamma-1)}.
\]
The choices of \(\gamma\), \(P\), \(j\), and \(k\) were arbitrary in their stated ranges, so the bound holds uniformly as claimed.
% lean: CausalSmith.Stat.WeakOverlap.orderedTreatedMass_lower
\end{enumerate}
\end{proof}
